\section{RNN, Convolution y Attention: la capacidad y la entrenabilidad son cosas distintas}

La formulation general de la sección anterior todavía depende del backbone. El curso no presenta a LSTM como un modelo fallido; al contrario, enfatiza que la recurrent network es muy expressive, al grado de poder aprender máquinas de estados semejantes a un contador. La motivación del Transformer proviene principalmente del computation graph: la dependency serial en secuencias largas, la path length entre posiciones distantes y el aprovechamiento del hardware limitaban la escala del entrenamiento.

\subsection{LSTM es muy Expressive, pero las actualizaciones de estado deben esperar su turno}

LSTM (Long Short-Term Memory) mitiga, mediante un gated cell state, el problema de los gradientes de largo plazo de la RNN convencional. Puede aprender transition complejas y, tanto en la teoría como en los experimentos, es capaz de simular ciertos comportamientos de conteo, copia y condicionales; sin embargo, cada paso temporal depende del resultado del paso anterior, de modo que, por muy capaz que sea el modelo, no puede calcular simultáneamente todos los hidden states de la secuencia.

\lecturefigure{slide-06-lstm-expressive.jpg}{LSTM puede implementar conteo y transformaciones de secuencias complejas}{Diapositiva recuperada de la grabación oficial V006; 00:12:09}

\begin{knowledgebox}{Lectura de la figura: no confundir un cuello de botella del sistema con falta de capacidad expresiva}
La tarea de la izquierda muestra que el recurrent state puede representar regularidades estructuradas de una secuencia; el diagrama de compuertas de la derecha indica que la actualización del estado cuenta con varias rutas aprendibles. Primero hay que ver qué puede expresar y después cómo debe programarse el cálculo; la figura respalda que LSTM es potente, pero no cambia el critical path en el que $h_t$ depende de $h_{t-1}$.
\end{knowledgebox}

\lecturefigure{slide-07-sequential-bottleneck.jpg}{El cuello de botella serial de recurrent reading y writing}{Diapositiva recuperada de la grabación oficial V007; 00:12:33}

\begin{warningbox}{Un lote paralelo no equivale a una secuencia paralela}
La GPU puede procesar al mismo tiempo muchas oraciones y muchas hidden dimension, pero aun así debe avanzar la recurrence de una misma oración en el orden $t=1,2,\ldots,n$. Aumentar el tamaño de lote no elimina el sequential critical path de cada muestra; las secuencias largas reducen el aprovechamiento del hardware y alargan la ruta de propagación del gradiente.
\end{warningbox}

\teachervoice{El ponente dice expresamente que las LSTM son «fantastic models» y, con el ejemplo $a^n\mapsto b^n$, muestra que una sola cell puede formar un contador aproximado. Lo que la clase objeta en realidad es atar esta gran capacidad expresiva a una ejecución serial token por token, no negar la recurrence en sí misma.}

\subsection{Convolution: ya es paralela, pero las relaciones a larga distancia dependen de apilar capas}

La convolution unidimensional puede procesar todas las posiciones al mismo tiempo, y cada capa agrega información dentro de una ventana local. El \term{receptive field (campo receptivo)} es el rango de la entrada del que puede depender una posición de salida; para una convolución de $L$ capas con kernel width $k$ y sin dilation, el campo receptivo es aproximadamente:
\[
R_L=1+L(k-1).
\]
donde $R_L$ es el rango visible en la capa $L$ y $k$ es el kernel width. Si $k=3$, cada capa adicional solo se extiende una posición hacia cada lado, por lo que una dependencia que abarque una longitud $n$ todavía requiere $O(n)$ capas o diseños adicionales como dilation/sparsity.

\lecturefigure{slide-08-convolution-parallel-structure.jpg}{Un convolutional sequence model permite calcular las posiciones en paralelo}{Diapositiva recuperada de la grabación oficial V008; 00:14:09}

En esta diapositiva, primero hay que ver si todas las posiciones pueden entrar simultáneamente a la misma capa y luego cuántos vecinos cubre cada kernel. La convolución elimina la cadena temporal recurrent y se adapta bien a la GPU; pero cada capa solo agrega información local. La figura respalda la mejora en parallelization, pero no muestra que dos tokens cualesquiera interactúen directamente en una sola capa.

\lecturefigure{slide-09-convolution-receptive-field.jpg}{El receptive field de la convolución crece gradualmente con el número de capas}{Diapositiva recuperada de la grabación oficial V009; 00:14:42}

\begin{knowledgebox}{Lectura de la figura: comparar la ruta más corta, no solo contar FLOPs}
Las rutas azul y roja representan la propagación de la información entre capas. Primero hay que determinar cuántas capas debe atravesar la información entre dos tokens distantes y luego compararlo con la conexión de un solo salto de la self-attention; un path más largo aumenta la dificultad de optimización y la pérdida de información. La dilated convolution puede acelerar la expansión, pero sigue exigiendo que el patrón de conexiones se defina de antemano, y no puede elegir vecinos dinámicamente según el contenido.
\end{knowledgebox}

\subsection{Encoder--Decoder Attention: leer la Source según su contenido}

La \term{encoder--decoder attention (atención cruzada)} permite que la target position asigne pesos a todos los source states según la decoder query actual. Supera el bottleneck del vector fijo y convierte la word alignment en un cálculo diferenciable. Si el encoder state es $h_j$ y el decoder state es $s_t$, puede escribirse como:
\[
e_{tj}=a(s_{t-1},h_j),\qquad
\alpha_{tj}=\frac{\exp e_{tj}}{\sum_{k}\exp e_{tk}},\qquad
c_t=\sum_j\alpha_{tj}h_j.
\]
donde $e_{tj}$ es el compatibility score, $\alpha_{tj}$ es la alignment normalizada y $c_t$ es el source context que lee el target step $t$. La attention ya accede a la source en paralelo, pero la estructura principal del encoder y del decoder todavía puede ser recurrent.

\lecturefigure{slide-10-encoder-decoder-attention.jpg}{La encoder--decoder attention establece conexiones por contenido entre source y target}{Diapositiva recuperada de la grabación oficial V010; 00:15:54}

\begin{knowledgebox}{Lectura de la figura: las flechas no son conexiones fijas, sino pesos que dependen de los datos}
La posición objetivo $y_t$ puede conectarse con todos los source states $x_j$; el grosor o la selección lo determina el contenido de la query y la key. Primero hay que ver en qué se distingue del offset fijo de la convolución y después notar que el decoder sigue generando en orden temporal; la figura respalda la content-based retrieval, pero todavía no reemplaza con self-attention el aprendizaje de representaciones de la propia source sentence.
\end{knowledgebox}

\teachervoice{Vaswani remonta la intuición de la attention a non-local means y a texture synthesis: extraer información de patches lejanos según la similitud de contenido. Esto muestra que «content talks to content» es un motif que surge de forma natural en distintas modalidades; lo esencial del Transformer es convertirlo en la estructura principal, no demostrar que el NLP posee un mecanismo misterioso exclusivo.}

\subsection{Resumen de la sección}

Las RNN aportan una gran capacidad expresiva, pero son seriales; la convolution aporta paralelismo, pero el path entre posiciones distantes sigue creciendo con la profundidad; la encoder--decoder attention, por su parte, demuestra que el content-based global lookup es eficaz. El siguiente paso de la self-attention es que cada token construya directamente una nueva representación a partir de todos los tokens de la misma secuencia.
